\unitopen{11}
\begin{candos}
\item poprosić o pomoc, o pozwolenie i o powtórzenie (\de{Kannst du \dots?}, \de{Darf ich \dots?});
\item powiedzieć, co musisz, możesz, chcesz, wolno Ci i powinieneś zrobić;
\item zbudować zdanie z \emph{klamrą}: czasownik modalny na 2.\ miejscu, bezokolicznik na końcu;
\item odróżnić \de{nicht müssen} od \de{nicht dürfen}.
\end{candos}

\section{Słownictwo}
\vocabtable{11}

\section{Odmiana}
\begin{center}\small
\begin{tabular}{@{}llllll@{}}
\toprule
 & \de{können} & \de{müssen} & \de{wollen} & \de{dürfen} & \de{sollen}\\
\midrule
\de{ich} & \de{kann} & \de{muss} & \de{will} & \de{darf} & \de{soll}\\
\de{du} & \de{kannst} & \de{musst} & \de{willst} & \de{darfst} & \de{sollst}\\
\de{er/sie/es} & \de{kann} & \de{muss} & \de{will} & \de{darf} & \de{soll}\\
\de{wir} & \de{können} & \de{müssen} & \de{wollen} & \de{dürfen} & \de{sollen}\\
\de{ihr} & \de{könnt} & \de{müsst} & \de{wollt} & \de{dürft} & \de{sollt}\\
\de{sie/Sie} & \de{können} & \de{müssen} & \de{wollen} & \de{dürfen} & \de{sollen}\\
\bottomrule
\end{tabular}
\end{center}
Pierwsza i trzecia osoba liczby pojedynczej \emph{nie mają końcówki} (\de{ich kann}, \de{er kann}). Z \de{möchte} z jednostki 8
to jest szósty czasownik modalny.

\section{Klamra zdaniowa}
Czasownik modalny (odmieniony) stoi na \emph{drugiej} pozycji, a znaczący czasownik w bezokoliczniku na \emph{końcu}:

\begin{center}\small
\begin{tabular}{@{}lllll@{}}
\de{Ich} & \textbf{\de{kann}} & \de{heute nicht ins Büro} & \textbf{\de{kommen}}.\\
\de{Morgen} & \textbf{\de{muss}} & \de{ich früh} & \textbf{\de{aufstehen}}.\\
\de{Können} & \textbf{\de{Sie}} & \de{das bitte} & \textbf{\de{wiederholen}}?\\
\end{tabular}
\end{center}
Gdy kontekst jest oczywisty, bezokolicznik można pominąć: \de{Ich muss los.} \de{Ich kann Deutsch.} \de{Das darfst du nicht.}

\begin{zasada}[Znaczenia]
\de{können} --- zdolność lub możliwość (\de{Ich kann Python.}). \de{müssen} --- konieczność. \de{wollen} --- silna wola
(\de{Ich will das lernen.}). \de{dürfen} --- pozwolenie (\de{Darf ich?}). \de{sollen} --- zobowiązanie od kogoś
(\de{Ich soll dich anrufen.}). \de{möchten} --- grzeczne życzenie.
\end{zasada}

\begin{falle}
\textbf{\de{nicht müssen} $\neq$ \de{nicht dürfen}.} \de{Du musst nicht kommen} znaczy „nie musisz” (możesz, ale nie
trzeba), a \de{Du darfst nicht kommen} znaczy „nie wolno Ci”. Polskie „nie musisz” i „nie wolno” są rozdzielone,
ale angielskie \emph{mustn't} / \emph{don't have to} tak samo, więc dodatkowe ryzyko dotyczy czytania
komunikatów: \de{Hier darf man nicht parken.} Nie ma tu „nie musisz”. Druga: \de{Ich muss gehen}, nie \de{Ich muss zu gehen}
(brak \de{zu} po modalnych).
\end{falle}

\begin{szwajcaria}
\de{Können Sie bitte \dots} jest w Szwajcarii bardziej użyteczne niż rozkaz; w pracy ludzie często mówią
\de{Könntest du kurz \dots}. Dialektalne \de{gäll} i \de{oder?} na końcu zdania („prawda?”) usłyszysz często;
w standardzie odpowiadasz po prostu \de{ja}, \de{genau}, \de{stimmt}.
\end{szwajcaria}

\begin{praca}
\de{Kannst du bitte den Pull Request reviewen?} \de{Ich muss das noch mit dem Team klären.}
\de{Wir sollen bis Freitag fertig sein.} \de{Darf ich dir kurz etwas zeigen?} \de{Können Sie das
bitte noch einmal erklären?} Zwroty, które ratują rozmowę: \de{Das habe ich nicht verstanden}, \de{Können Sie das anders sagen?}
\end{praca}

\section{Dialogi}
\textbf{Dialog 1 --- prośba o pomoc w pracy}
\begin{dialog}
\turn{A}{Entschuldigung, kannst du mir kurz helfen?}{Przepraszam, możesz mi na chwilę pomóc?}
\turn{B}{Natürlich. Was kann ich für dich tun?}{Oczywiście. Co mogę dla ciebie zrobić?}
\turn{A}{Ich muss den Server neu starten, aber ich darf das nicht allein machen.}{Muszę zrestartować serwer, ale nie wolno mi tego robić samemu.}
\turn{B}{Ich komme gleich. Wir müssen erst den Chef fragen.}{Zaraz przyjdę. Najpierw musimy zapytać szefa.}
\turn{A}{Danke. Kannst du das auch dokumentieren?}{Dzięki. Możesz to też udokumentować?}
\end{dialog}

\section{Ćwiczenia}
\begin{cwiczenia}
\exitem{Uzupełnij: \de{Ich \dots{} Deutsch lernen.} (wollen)}{\de{will}}
\exitem{Uzupełnij: \de{Du \dots{} heute arbeiten.} (müssen)}{\de{musst}}
\exitem{Uzupełnij: \de{Er \dots{} gut programmieren.} (können)}{\de{kann}}
\exitem{Uzupełnij: \de{Wir \dots{} hier nicht rauchen.} (dürfen)}{\de{dürfen}}
\exitem{Uzupełnij: \de{Ihr \dots{} um zehn Uhr da sein.} (sollen)}{\de{sollt}}
\exitem{Ułóż zdanie: \de{Ich / heute / kommen / nicht / kann / ins Büro}}{\de{Ich kann heute nicht ins Büro kommen.}}
\exitem{Ułóż zdanie, zaczynając od \de{Morgen}: \de{muss / ich / früh / aufstehen}}{\de{Morgen muss ich früh aufstehen.}}
\exitem{Zadaj pytanie: czy możesz powtórzyć (formalnie).}{\de{Können Sie das wiederholen?}}
\exitem{Zapytaj o pozwolenie: czy mogę tu zaparkować?}{\de{Darf ich hier parken?}}
\exitem{Co znaczy \de{Du musst nicht kommen}?}{Nie musisz przyjść (możesz, ale nie trzeba).}
\exitem{Co znaczy \de{Du darfst nicht kommen}?}{Nie wolno Ci przyjść.}
\exitem{Popraw: \de{Ich muss zu gehen.}}{\de{Ich muss gehen.}}
\exitem{Przetłumacz: „Muszę to jeszcze wyjaśnić z zespołem.”}{\de{Ich muss das noch mit dem Team klären.}}
\end{cwiczenia}

\section{Trening słownictwa (generowany)}
\drill{11}

\weekend{Napisz osiem zdań z różnymi czasownikami modalnymi o swojej pracy i nauce
(co musisz, co możesz, co chcesz, czego Ci nie wolno). Przekształć trzy z nich w prośby do kolegi (\de{Kannst du \dots?}).
Zagraj z AI scenkę „prośba o pomoc”.}
